% 问题四：总体路线

本问题在问题 2 的消解框架基础上，为 C 类装备新增"调整间隔时长"这一动作维度，扩展了动作空间。

\textbf{动作空间扩展}

对于每个计划 $P_i$，可选动作集合为：
\begin{itemize}
    \item \textbf{保留}：不调整任何参数
    \item \textbf{频段平移}：$\Delta f \in [-10, 10]$，满足 $f_s + \Delta f \geq 0$ 且 $f_e + \Delta f \leq 100$
    \item \textbf{时间平移}：$\Delta t \in [-5, 5]$，满足 $t_s + \Delta t \geq 0$
    \item \textbf{间隔调整}（仅 C 类）：$\Delta \delta \in [-10, 10]$，满足 $\delta + \Delta \delta \geq 2$（避免自干扰）
    \item \textbf{撤销}
\end{itemize}

每个计划仍然\textbf{最多只调整一个参数}，因此间隔调整与频段/时间平移互斥。

\textbf{三种求解方法}

\begin{enumerate}
    \item \textbf{贪心 + 优先级 + 间隔维度}：按 A > B > C 优先级排序，同类内按冲突度数降序。对每个计划依次尝试：保留 → 频段平移（$\pm 1$ 到 $\pm 10$）→ 时间平移（$\pm 1$ 到 $\pm 5$）→ 间隔调整（仅 C 类，$\pm 1$ 到 $\pm 10$）→ 撤销。每档内按幅度从小到大枚举。

    \item \textbf{模拟退火 + 间隔维度}：以贪心解为初始解，随机扰动时可选择更换为间隔调整动作。使用多种子策略，退火结束后收尾消毒保证 0 冲突。

    \item \textbf{贪心 + 局部搜索 + 间隔维度}：在贪心解基础上进行两阶段局部改进——撤销恢复（逐个尝试把被撤销计划重新安置，允许使用间隔调整动作）与调整瘦身（把已调整计划回退为保留或更小幅度的动作），探索间隔调整的改进空间。

    \item \textbf{整数线性规划 ILP + 间隔动作}：将间隔调整作为新的动作类型纳入 ILP 模型，以撤销数最小化为目标函数。
\end{enumerate}

\textbf{最优方案选择}：比较四种方法的消解结果，按分层优先级（撤销数 → A 撤销 → 调整数 → 幅度）选择最优方案。同时与问题 2 的最优结果（模拟退火，撤销 50 个）对比，评估放宽间隔约束带来的收益。
